\printedpage{29}
% head: A. Grothendieck : Dérivateurs, Ch. VI, Édition des Universités de Montpellier II et Paris VII
% foot: Version 9 novembre 2001

\begin{enumerate}[leftmargin=*]
\item[d)] For $f : X \longrightarrow Y$ to be in $W$, it is necessary and sufficient that $\pi_0(f)$ be bijective, and that the $f_i : X_i \longrightarrow Y_i$ induced on the connected components be in $W$.
\end{enumerate}

These axioms are satisfied by all the $W$ that I know, and that one would feel like regarding as localizers. On the other hand, L5, L6, L6 bis (p. 43, 44) are axioms of a much more special nature. I feel like saying that they distinguish localizers of a strictly \textit{``homotopical''} nature from those of a \textit{``homological''} (or ``cohomological'') nature. I should make sure that they are satisfied for the $W_i$ ($i$-homotopy equivalence) (\footnote[32]{This is \textit{false}, cf. annotation above.}). % [sic] qui je connais -> que je connais, relative pronoun object requires "que"

The fact that L4 bis, i.e. c) above, implies the dual axiom, follows from the fact that L4 bis implies this

\textbf{Proposition 2.} \textit{Let $f : X \longrightarrow Y$ be in} $\Cat$\textit{, then $f \in W$ if and only if $f^o : X^o \longrightarrow Y^o$ is in $W$.}

Seen by the argument of I, page 75 (cor. 2) - uses only L4, not even L1, L3.

\msp{page 46}

I am going to try to summarize what I think I have understood:

\begin{enumerate}[leftmargin=*]
\item[1)] $W$-equivalence ``by fibers'' is of interest (for constructing localized categories) only in the case where one knows that it already implies, either local $W$-equivalence, or colocal $W$-equivalence (or even universal $W$-equivalence) over $S$, hence when one assumes $X$, $Y$ to be at the very least both either $W$-smooth, or $W$-proper, or $W$-fibered over $S$.

\item[2)] One could call local $W$-equivalence over $S$ \textit{``left fine $W$-equivalence''} over $S$, and colocal $W$-equivalence over $S$ \textit{``right fine $W$-equivalence''} over $S$. They coincide in $\uc{Fibhot}\,S$, a fortiori in $\uc{Parf}\,S$ and in the four categories (cf. diagram p. 8) $\uc{Fib}_0S$, $\uc{Fibsc}_0S$, $\uc{Cofib}_0S$, $\uc{Cofibsc}_0S$ (\footnote[33]{One could call them the six \textit{inner} categories of the diagram.}). On the other hand, they do not coincide on the \textit{bordering} categories of the diagram, in particular on $\uc{Liss}\,S$, $\uc{Prop}\,S$, $\uc{Fib}\,S$, $\uc{Cofib}\,S$, $\uc{Fibsc}\,S$, $\uc{Cofibsc}\,S$. In this regard, it is the \textit{fine $W$-equivalence} \textbf{on the right} (in terms of the $s\backslash X$) that is of interest on the (alas) left side of the diagram, i.e. for $\uc{Liss}\,S$, $\uc{Fib}\,S$, $\uc{Fibsc}\,S$, whereas it is the \textit{fine $W$-equivalence} \textbf{on the left} (in terms of the $X/s$) that is relevant on the right side, for $\uc{Prop}\,S$, $\uc{Cofib}\,S$, $\uc{Cofibsc}\,S$. One will never have to look simultaneously at the two fine equivalences on one and the same category of the diagram, with the sole exception of $\Cat/S$ as a whole.
\end{enumerate}

\msp{page 47}

\begin{enumerate}[leftmargin=*]
\item[3)] Coarse $W$-equivalence, for its part, is ``good everywhere'' - but that depends for whom. For the 6 categories that are above $\uc{Fibhot}\,S$, i.e.
\[
\uc{Fibhot}\,S,\quad \uc{Parf}\,S,\quad \uc{Fib}_0S,\quad \uc{Fibsc}_0S,\quad \uc{Cofib}_0S,\quad \uc{Cofibsc}_0S,
\]
\printedpage{30}
% head: A. Grothendieck : Dérivateurs, Ch. VI, Édition des Universités de Montpellier II et Paris VII
% foot: Version 9 novembre 2001
it coincides with all the other equivalences (assuming axiom L5), for the three upper left bordering categories % [sic] moyennnant -> moyennant, triple n
\[
\uc{Liss}\,S,\quad \uc{Fib}\,S,\quad \uc{Fibsc}\,S
\]
of the diagram, it coincides with right fine $W$-equivalence, and dually on
\[
\uc{Prop}\,S,\quad \uc{Cofib}\,S,\quad \uc{Cofibsc}\,S
\]
it coincides with left fine $W$-equivalence. As for $\Cat/S$ as a whole, it is the only category of the diagram where there are \textit{three} different interesting notions of $W$-equivalence, namely: the coarse one, and the two fine ones (left, right). On the other six bordering categories of the diagram, there are in principle two: the coarse one, and a specified one of the two fine ones (forbidden to look at the other - at one's own risk and peril \ldots): the right one on the left side, the left one on the right side. (NB I ought to redo the diagram, really, putting the fibrations on the right, not on the left.) As for the six inner categories of the diagram, in them there is only a single notion of $W$-equivalence left (since on them the four notions coincide).
\end{enumerate}

\msp{page 48 (vide)}

\msp{page 49} (\textbf{NB} there is no page 48)

I now summarize what I think I have understood about the corresponding localized categories. There are, up to equivalence, only \textit{three} different ones all told, when one takes the ``admissible'' localizers in the categories, as I have just said (that being in short the objective interest of the prescriptions!) To be precise:

\begin{enumerate}[leftmargin=*]
\item[$1^o$)] The \textit{coarse $W$-equivalence} gives ``the same'' localized category for the 13 categories of the diagram, i.e. all the transition arrows of the diagram induce equivalences for this coarse $W$-equivalence. The localized category will be denoted \fbox{$\HOT^{\mathrm{lc}}_W(S)$}, where the exponent lc means: (relative $W$-homotopy types over $S$) \textit{``locally constant.''} (I have only a small doubt left for $\uc{Fibsc}\,S$ and $\uc{Cofibsc(S)}$ (\footnote[34]{This doubt is unfounded, everything works perfectly!}), whether coarse equivalence there is really and truly reasonable. We will see when doing the proof.)

\item[$2^o$)] The \textit{right fine $W$-equivalence} gives the same result for all the categories of the diagram where it is considered as admitted, \textit{and where it does not coincide with the coarse $W$-equivalence} already considered, hence on the four left bordering categories
\[
\Cat/S,\quad \uc{Liss}\,S,\quad \uc{Fib}\,S,\quad \uc{Fibsc}\,S
\]
(\footnote[35]{\textbf{NB} On the last three this equivalence is $W$-equivalence by fibers.}). The localized category is denoted \fbox{$\HOT_W(S)$}.
\end{enumerate}

\msp{page 50}

\printedpage{31}
% head: A. Grothendieck : Dérivateurs, Ch. VI, Édition des Universités de Montpellier II et Paris VII
% foot: Version 9 novembre 2001

\begin{enumerate}[leftmargin=*]
\item[$2'^o$)] Dually, the \textit{left fine $W$-equivalence} gives the same result on all the categories of the diagram where it is admitted, and where it does not coincide with the coarse one, i.e. on the 4 categories
\[
\Cat/S,\quad \uc{Prop}\,S,\quad \uc{Cofib}\,S,\quad \uc{Cofibsc}\,S.
\]
It is appropriate to denote the common category of fractions by the symbol
\[
\HOT_W(S^o)^o \qquad \text{(or also $\HOT^o_W(S)$)},
\]
since it is isomorphic to the one denoted by this symbol.
\end{enumerate}

If one wants unambiguous notations, designating well-defined categories, and not just ``up to equivalence'', I will choose
\[
\HOT_W(S) \defeq \uc{Fibsc}(S)(W^{\mathrm{d}}_S)^{-1}
\]
whence
\[
\HOT_W(S^o) \qquad \underbrace{\simeq}_{\text{\footnotesize isomorphism of categories}} \qquad \uc{Cofibsc}(S)(W^{\mathrm{g}}_S)^{-1},
\]
where $W^{\mathrm{d}}_S$, $W^{\mathrm{g}}_S$ denote the sets of arrows corresponding to right resp. left fine $W$-equivalence. For $\HOT^{\mathrm{lc}}_W(S)$, I prefer a definition as symmetric as possible in $S$, $S^o$, and moreover a category as small as possible, which leads me to choose
\[
\HOT^{\mathrm{lc}}_W(S) \defeq (\uc{Parf}\,S)W^{-1}_S,
\]
where we have denoted by $W_S$ the $W$-equivalences in $\uc{Parf}\,S$. (No need this time to specify which ones - all four \texttt{[notions]} coincide in $\uc{Parf}\,S$!)

\msp{page 51}

Let us note the isomorphism
\[
\uc{Parf}(S^o) \simeq \uc{Parf}\,S
\]
(coming from the fact that $f : X \longrightarrow Y$ in $\Cat$ is a perfect fibration if and only if $f^o$ is). This isomorphism $X \longmapsto X^o$ is compatible with the localizers (since $f \in W \Longleftrightarrow f^o \in W$), whence
\[
\HOT^{\mathrm{lc}}_W(S^o) \simeq \HOT^{\mathrm{lc}}_W(S)
\]
(isomorphism of categories). We will have two canonical functors between the three categories of relative $W$-homotopy types over $S$
\[
\begin{tikzcd}[column sep=small]
(\uc{Fibsc}_0S)W_S^{-1} \arrow[rr, "\approx"] \arrow[dr, dotted, "i'"'] & & \HOT^{\mathrm{lc}}_W(S) \arrow[dl, "i"] \arrow[dr, "j"'] & & (\uc{Cofibsc}_0S)W_S^{-1} \arrow[ll, "\approx"'] \arrow[dl, dotted, "j'"] \\
& \HOT_W(S) & & \HOT^o_W(S) \defeq \HOT_W(S^o). &
\end{tikzcd}
\]
\printedpage{32}
% head: A. Grothendieck : Dérivateurs, Ch. VI, Édition des Universités de Montpellier II et Paris VII
% foot: Version 9 novembre 2001
\texttt{[rather} $\defeq \HOT_W(S^o)^o$\texttt{]} These arrows $i$, $j$ are suggested by the dotted arrows $i'$, $j'$, coming from the inclusions % [sic] suggerées -> suggérées, missing accent
\[
\uc{Fibsc}_0S \lhook\joinrel\longrightarrow \uc{Fibsc}\,S,\quad \uc{Cofibsc}_0S \lhook\joinrel\longrightarrow \uc{Cofibsc}\,S,
\]
which would define $i$, $j$ with the help of a choice of quasi-inverse functors for the horizontal arrows. But $i$, $j$ would be defined only up to (unique) isomorphism. It will be prettier to choose the quasi-inverses in a canonical way, with the help of the arrows
\[
\uc{Fibsc}_0S \xleftarrow{\Phi_0} \uc{Parf}\,S \xrightarrow{\Psi_0} \uc{Cofibsc}_0S
\]
(where $\Phi_0 = \Phi|\uc{Parf}\,S$, $\Psi_0 = \Psi|\uc{Parf}\,S$) of the diagram p. 8. The main question
\msp{page 52}
which remains unelucidated, is whether these functors $i$, $j$ really are \textit{fully faithful,} as I certainly expect.

I now want to prove the preceding assertions about the localized categories.

\textcircled{1} The localized category for coarse $W$-equivalence. It will suffice, up to an $\varepsilon$, to establish the conditions of the lemma p. 11, in the following three cases.
\begin{tikzcd}[baseline=(\tikzcdmatrixname-1-1.base)]
\uc{Parf}\,S \arrow[r, "r"] & \Cat/S \arrow[l, dotted, bend left, "C"]
\end{tikzcd}
(that will give the case of $\uc{Fibhot}\,S$, $\uc{Liss}\,S$ and $\uc{Prop}\,S$ over $\Cat/S$) and
\begin{tikzcd}[baseline=(\tikzcdmatrixname-1-1.base)]
\uc{Fibsc}_0S \arrow[r, "p"] & \Cat/S \arrow[l, dotted, bend left]
\end{tikzcd},
\begin{tikzcd}[baseline=(\tikzcdmatrixname-1-1.base)]
\uc{Cofibsc}_0S \arrow[r, "q"] & \Cat/S \arrow[l, dotted, bend left]
\end{tikzcd},
which will give respectively the cases of
\[
\uc{Fib}_0S,\ \uc{Fib}\,S,\ \uc{Liss}\,S \qquad \text{and} \qquad \uc{Cofib}_0S,\ \uc{Cofib}\,S,\ \uc{Prop}\,S \qquad \text{over } S.
\]
But the trouble for $p$ and $q$, is that the functors in the reverse direction that I had constructed (and again denoted $\Phi_0$, $\Psi_0$ (but they were essentially $\Phi\circ r\circ C$ and $\Psi\circ r\circ c$, not to be confused with the functors $\Phi_0$, $\Psi_0$ at the bottom of the preceding page, and also marked in the diagram p. 8), let us rather denote them $\Phi'$, $\Psi'$, work well on the $\Cat/S$ side, i.e. $p\Phi'$, $q\Psi'$ receive the functor $\mathrm{id}_{\Cat/S}$ by a $W$-arrow, but not on the side of the target categories - one cannot manage to find a natural arrow between $\Phi'p$ resp. $\Psi'q$, and the identity functors of $\uc{Fibsc}_0S$ resp. $\uc{Cofibsc}_0S$. I do
\msp{page 53}
not see, indeed, that there are any - the difficulty, by this route, seems irremediable. On the other hand, if one regards $\Phi'$, $\Psi'$ as functors with values in $\uc{Fib}_0S$, $\uc{Cofib}_0S$ respectively (thus in the categories where one does not require the arrows to be cartesian resp. cocartesian), one again finds $W$-morphisms % [sic] catésiennes -> cartésiennes, missing r
\[
\mathrm{id} \longrightarrow \Phi'p, \qquad \mathrm{id} \longrightarrow \psi'q, % [sic] \psi'q -> \Psi'q, lowercase psi where \Psi' is used everywhere else
\]
which at least already settles the case of all the categories of the diagrams p. 8 that are
\printedpage{33}
% head: A. Grothendieck : Dérivateurs, Ch. VI, Édition des Universités de Montpellier II et Paris VII
% foot: Version 9 novembre 2001
full subcategories of $\Cat/S$ (\footnote[36]{But it is idiotic to take such complicated functors, because one wants to reduce everything to $\Cat/S$. It is moreover simple to compare $\uc{Fibhot}\,S$ directly with $\uc{Fibsc}_0S$ and with $\uc{Cofibsc}_0S$. $\Phi_0$ and $\Psi_0$ (which are in fact defined on $\uc{Fibhot}$, not just on $\uc{Parf}$), and one finds as below $\mathrm{id} \xrightarrow{W} \pi_0\Phi_0$ and $\Phi_0\pi_0 \xrightarrow{W} \mathrm{id}$ (likewise for $\Psi_0$ and $\uc{Cofibsc}_0S$), where $\pi_0 : \uc{Fibsc}_0S \longrightarrow \uc{Fibhot}\,S$.}). It remains only to treat the case of the four categories of ``split'' structures
\[
\uc{Fibsc}\,S,\quad \uc{Fibsc}_0S \qquad \text{and} \qquad \uc{Cofibsc}\,S,\quad \uc{Cofibsc}_0S.
\]
Of course it suffices to do it on one side, say on the fibered side, by duality ($S \longmapsto S^o$) this will establish the symmetric result. So we are going to look a little more closely at the square of functors
\[
\begin{tikzcd}[column sep=huge, row sep=large]
\uc{Fibsc}\,S \arrow[d, "p"] & \uc{Fibsc}_0S \arrow[l, hook, "\alpha'"'] \arrow[d, "p_0"'] \\
\uc{Fib}\,S \arrow[u, dotted, bend left, "\varphi"] & \uc{Fib}_0S \arrow[l, hook, "\alpha"'] \arrow[u, dotted, bend right, "\varphi_0"']
\end{tikzcd}
\]
(\footnote[37]{NB It is a matter of proving only that $p$, $p_0$ induce equivalences for the localized categories (for coarse $W$-equivalence). We already know that this is so for $\alpha$, and it will follow for $\alpha'$.}, \footnote[38]{needless to look at this square, cf. marginal annotations above (top of the page)}). The functor
\[
\Phi : \Cat/S \longrightarrow \uc{Fibsc}\,S
\]
induces functors in the reverse direction to $p$, $p_0$ on the full subcategories $\uc{Fib}\,S$, $\uc{Fib}_0S$ of $\Cat/S$, functors marked by dotted lines as $\varphi$, $\varphi_0$.

\msp{page 54}

It is only necessary to check that if $X/S$ is in $\uc{Fib}_0S$, then $\varphi(X) = \Phi(X)$ is in the full subcategory $\uc{Fibsc}_0S$ of $\uc{Fibsc}\,S$. But we have seen that this is the case, provided only that $X/S$ is in $\uc{Fibhot}\,S$, for the base change morphisms in $\Phi(X)$, associated to the arrows $u : s \longrightarrow t$ in $S$
\[
u^*_{\Phi_X} : t\backslash X \longrightarrow s\backslash X
\]
are $W_\infty$-equivalences if $X$ is Hot-fibered over $S$. On the other hand, for arbitrary $X$ in $\Cat/S$, one has an arrow, functorial in $X$,
\[
X \longrightarrow \Phi(X),
\]
i.e. $\mathrm{id}_{\Cat/S} \longrightarrow \pi\Phi$ where $\pi : \uc{Fibsc}\,S \longrightarrow \Cat/S$ is the canonical functor, which is a $W_\infty$-equivalence (even a homotopism), hence a $W$-equivalence. A fortiori the functors $p\varphi$, $p_0\varphi_0$ give rise ($\uc{Fib}\,S$ and $\uc{Fib}_0S$ being full subcategories of $\Cat/S$) to arrows
\[
\mathrm{id} \longrightarrow p\varphi, \qquad \mathrm{id} \longrightarrow p_0\varphi_0.
\]